\begin{center}Michael Renardy\quad Robert C. Rogers\\An Introduction to Partial Differential Equations\\Second Edition\\With 41 Illustrations\\Springer\end{center}

\clearpage\section*{Series Preface}\phantomsection\label{sec:series-preface}\addcontentsline{toc}{section}{Series Preface}

Mathematics is playing an ever more important role in the physical and biological sciences, provoking a blurring of boundaries between scientific disciplines and a resurgence of interest in the modern as well as the classical techniques of applied mathematics. This renewal of interest, both in research and teaching, has led to the establishment of the series Texts in Applied Mathematics (TAM).

The development of new courses is a natural consequence of a high level of excitement on the research frontier as newer techniques, such as numerical and symbolic computer systems, dynamical systems, and chaos, mix with and reinforce the traditional methods of applied mathematics. Thus, the purpose of this textbook series is to meet the current and future needs of these advances and to encourage the teaching of new courses.

TAM will publish textbooks suitable for use in advanced undergraduate and beginning graduate courses, and will complement the Applied Mathematical Sciences (AMS) series, which will focus on advanced textbooks and research-level monographs.

\begin{flushright}Pasadena, California\hfill J.E. Marsden\\Providence, Rhode Island\hfill L. Sirovich\\College Park, Maryland\hfill S.S. Antman\end{flushright}

\clearpage\section*{Preface}\phantomsection\label{sec:author-preface}\addcontentsline{toc}{section}{Preface}

Partial differential equations are fundamental to the modeling of natural phenomena; they arise in every field of science. Consequently, the desire to understand the solutions of these equations has always had a prominent place in the efforts of mathematicians; it has inspired such diverse fields as complex function theory, functional analysis and algebraic topology. Like algebra, topology and rational mechanics, partial differential equations are a core area of mathematics.

Unfortunately, in the standard graduate curriculum, the subject is seldom taught with the same thoroughness as, say, algebra or integration theory. The present book is aimed at rectifying this situation. The goal of this course was to provide the background which is necessary to initiate work on a Ph.D. thesis in PDEs. The level of the book is aimed at beginning graduate students. Prerequisites include a truly advanced calculus course and basic complex variables. Lebesgue integration is needed only in Chapter~\ref{chap:nonlinear-elliptic-equations}, and the necessary tools from functional analysis are developed within the course.

The book can be used to teach a variety of different courses. Here at Virginia Tech, we have used it to teach a four-semester sequence, but (more often) for shorter courses covering specific topics. Students with some undergraduate exposure to PDEs can probably skip Chapter~\ref{chap:introduction}. Chapters~\ref{chap:characteristics}--\ref{chap:maximum-principles} are essentially independent of the rest and can be omitted or postponed if the goal is to learn functional analytic methods as quickly as possible. Only the basic definitions at the beginning of Chapter~\ref{chap:characteristics}, the Weierstrass approximation theorem and the Arzela-Ascoli theorem are necessary for subsequent chapters. Chapters~\ref{chap:nonlinear-elliptic-equations}, \ref{chap:energy-methods-evolution-problems} and \ref{chap:semigroup-methods} are independent of each other (except that Chapter~\ref{chap:semigroup-methods} uses some definitions from the beginning of Chapter~\ref{chap:energy-methods-evolution-problems}) and can be covered in any order desired.

We would like to thank the many friends and colleagues who gave us suggestions, advice and support. In particular, we wish to thank Pavel Bochev, Guowei Huang, Wei Huang, Addison Jump, Kyehong Kang, Michael Keane, Hong-Chul Kim, Mark Mundt and Ken Mulzet for their help. Special thanks is due to Bill Hrusa, who read a good deal of the manuscript, some of it with great care and made a number of helpful suggestions for corrections and improvements.

\section*{Notes on the second edition}\phantomsection\label{sec:second-edition-notes}

We would like to thank the many readers of the first edition who provided comments and criticism. In writing the second edition we have, of course, taken the opportunity to make many corrections and small additions. We have also made the following more substantial changes.

\begin{itemize}
\item We have added new problems and tried to arrange the problems in each section with the easiest problems first.
\item We have added several new examples in the sections on distributions and elliptic systems.
\item The material on Sobolev spaces has been rearranged, expanded, and placed in a separate chapter. Basic definitions, examples, and theorems appear at the beginning while technical lemmas are put off until the end. New examples and problems have been added.
\item We have added a new section on nonlinear variational problems with ``Young-measure'' solutions.
\item We have added an expanded reference section.
\end{itemize}
